\documentclass[11pt,a4paper]{article}
\usepackage[utf8]{inputenc}
\usepackage[T1]{fontenc}
\usepackage{geometry}
\geometry{margin=1in}
\usepackage{hyperref}
\usepackage{enumitem}
\usepackage{booktabs}
\usepackage{longtable}
\usepackage{xcolor}
\usepackage{titlesec}
\usepackage{amsmath,amssymb}

\titleformat{\section}{\Large\bfseries\color{blue!70!black}}{}{0em}{}
\titleformat{\subsection}{\large\bfseries\color{blue!50!black}}{}{0em}{}
\titleformat{\subsubsection}{\normalsize\bfseries\color{blue!30!black}}{}{0em}{}

\title{\textbf{Replication Plan}\\[0.5em]
\large Effect of Single-Atom Platinum Doping and Facet Dependence on the Electronic Structure and Light Absorption of La$_2$Ti$_2$O$_7$: A DFT Study\\[0.3em]
\normalsize OSTI ID: 1981773}
\author{Auto-generated Replication Blueprint}
\date{\today}

\begin{document}
\maketitle
\tableofcontents
\newpage

%---------------------------------------------------------------------
\section{Paper Overview}
%---------------------------------------------------------------------
This paper employs density functional theory (DFT) to investigate the effect of single-atom Pt doping on the electronic structure and optical absorption of monoclinic La$_2$Ti$_2$O$_7$ (LTO), a photocatalyst for water splitting. The study examines multiple surface facets ((010), (001), (100)), computes surface energies, analyzes band structures and density of states (DOS) upon Pt substitution at Ti sites, performs Bader charge analysis, and calculates optical absorption spectra using the DFT+U or hybrid functional approach.

%---------------------------------------------------------------------
\section{Detailed Workflow}
%---------------------------------------------------------------------

\subsection{Phase 1: Structure Preparation}
\begin{enumerate}[label=\textbf{1.\arabic*}]
  \item \textbf{Obtain the La$_2$Ti$_2$O$_7$ crystal structure.}
    \begin{itemize}
      \item From Materials Project or ICSD (monoclinic, space group $P2_1$).
      \item Optimize the bulk unit cell with DFT.
    \end{itemize}
  \item \textbf{Construct surface slab models.}
    \begin{itemize}
      \item Cleave (010), (001), (100) surfaces.
      \item Add vacuum layer ($\geq 15$\,\AA).
      \item Ensure stoichiometric and symmetric slabs.
    \end{itemize}
  \item \textbf{Introduce Pt dopant.}
    \begin{itemize}
      \item Substitute one Ti atom with Pt at various surface and subsurface sites.
      \item Create multiple doping configurations for each facet.
    \end{itemize}
\end{enumerate}

\subsection{Phase 2: DFT Calculations}
\begin{enumerate}[label=\textbf{2.\arabic*}]
  \item \textbf{Geometry optimization.}
    \begin{itemize}
      \item Relax all structures using PBE functional with DFT+U (or HSE06 for electronic properties).
      \item Convergence: forces $< 0.01$\,eV/\AA, energy $< 10^{-5}$\,eV.
      \item VASP settings: PAW pseudopotentials, energy cutoff ($\sim$520\,eV), $k$-point mesh (Gamma-centered).
    \end{itemize}
  \item \textbf{Surface energy calculations.}
    \begin{equation*}
      \gamma = \frac{E_{\text{slab}} - n \cdot E_{\text{bulk}}}{2A}
    \end{equation*}
  \item \textbf{Dopant formation energy.}
    \begin{equation*}
      E_f = E_{\text{doped}} - E_{\text{pristine}} - \mu_{\text{Pt}} + \mu_{\text{Ti}}
    \end{equation*}
\end{enumerate}

\subsection{Phase 3: Electronic Structure Analysis}
\begin{enumerate}[label=\textbf{3.\arabic*}]
  \item \textbf{Band structure calculations} along high-symmetry paths.
  \item \textbf{Projected density of states (PDOS)} --- decompose by atom and orbital.
  \item \textbf{Bader charge analysis} to quantify charge transfer upon doping.
  \item \textbf{Charge density difference} plots ($\Delta \rho$) for Pt--O bonding visualization.
\end{enumerate}

\subsection{Phase 4: Optical Properties}
\begin{enumerate}[label=\textbf{4.\arabic*}]
  \item \textbf{Compute dielectric function} $\varepsilon(\omega) = \varepsilon_1(\omega) + i\varepsilon_2(\omega)$.
  \item \textbf{Derive absorption coefficient} $\alpha(\omega)$ from the dielectric function.
  \item \textbf{Compare pristine vs.\ Pt-doped} absorption spectra; quantify visible-light absorption enhancement.
\end{enumerate}

\subsection{Phase 5: Validation}
\begin{enumerate}[label=\textbf{5.\arabic*}]
  \item Reproduce band gap values for pristine and doped systems.
  \item Reproduce surface energy ranking across facets.
  \item Reproduce DOS and absorption spectra plots.
  \item Compare Bader charges to paper's values.
\end{enumerate}

%---------------------------------------------------------------------
\section{Datasets}
%---------------------------------------------------------------------

\begin{longtable}{p{4.5cm} p{3.5cm} p{6cm}}
\toprule
\textbf{Dataset / Source} & \textbf{Type} & \textbf{Location / Notes} \\
\midrule
\endfirsthead
\toprule
\textbf{Dataset / Source} & \textbf{Type} & \textbf{Location / Notes} \\
\midrule
\endhead
La$_2$Ti$_2$O$_7$ crystal structure & CIF file & Materials Project: \url{https://materialsproject.org/} (mp-3533 or similar) \\
\midrule
PAW pseudopotentials & VASP POTCARs & Distributed with VASP license \\
\midrule
Experimental band gap & Reference & Literature values for validation \\
\midrule
Supplementary information & Tables, structures & Accompanying the paper \\
\bottomrule
\end{longtable}

%---------------------------------------------------------------------
\section{Models, Tools, and Software}
%---------------------------------------------------------------------

\begin{longtable}{p{4.5cm} p{2.5cm} p{6.5cm}}
\toprule
\textbf{Tool / Library} & \textbf{Version} & \textbf{Purpose / URL} \\
\midrule
\endfirsthead
\toprule
\textbf{Tool / Library} & \textbf{Version} & \textbf{Purpose / URL} \\
\midrule
\endhead
VASP & 5.4+ or 6.x & DFT calculations (geometry optimization, band structure, optics) \newline \url{https://www.vasp.at/} \\
\midrule
Quantum ESPRESSO & 7.0+ & Open-source DFT alternative \newline \url{https://www.quantum-espresso.org/} \\
\midrule
VESTA & $\geq 3.5$ & Crystal structure visualization \newline \url{https://jp-minerals.org/vesta/} \\
\midrule
pymatgen & $\geq 2022.0$ & Structure manipulation, VASP I/O \newline \url{https://pymatgen.org/} \\
\midrule
ASE & $\geq 3.22$ & Atomic simulation environment \\
\midrule
Bader analysis code & latest & Henkelman group Bader charge analysis \newline \url{http://theory.cm.utexas.edu/henkelman/code/bader/} \\
\midrule
VASPKIT & latest & Post-processing VASP outputs \newline \url{https://vaspkit.com/} \\
\midrule
Matplotlib & $\geq 3.4$ & Plotting band structures, DOS, spectra \\
\midrule
sumo & latest & Publication-quality band structure / DOS plots \newline \url{https://github.com/SMTG-Bham/sumo} \\
\bottomrule
\end{longtable}

\subsection{Hardware Requirements}
\begin{itemize}
  \item \textbf{HPC cluster}: DFT surface slab calculations require $\sim$16--64 cores, $\geq$ 64\,GB RAM.
  \item \textbf{HSE06 calculations}: significantly more expensive; $\sim$100+ cores, days per system.
  \item \textbf{Storage}: $\sim$50\,GB per slab calculation.
\end{itemize}

\end{document}
